\documentclass[aspectratio=169]{beamer}
\usetheme{Madrid}
\usecolortheme{default}
\usepackage{graphicx}
\usepackage{booktabs}
\usepackage{hyperref}
\usepackage{tabularx}
\usepackage{array}

\title{Oxford Pets Image Classification}
\subtitle{From Scratch CNN vs Transfer Learning}
\institute[THD]{6th Semester, B.Sc. Artificial Intelligence\\Technische Hochschule Deggendorf\\Professor: Prof. Dr. Isabel Hübener}
\date{Deep Learning and Big Data, Summer Semester 2026}

\setbeamertemplate{navigation symbols}{}
\setbeamertemplate{footline}{
  \leavevmode%
  \hbox{%
    \begin{beamercolorbox}[wd=.78\paperwidth,ht=2.25ex,dp=1ex,center]{title in head/foot}%
      \insertshorttitle
    \end{beamercolorbox}%
    \begin{beamercolorbox}[wd=.22\paperwidth,ht=2.25ex,dp=1ex,right]{date in head/foot}%
      \insertframenumber{} / \inserttotalframenumber\hspace*{2ex}
    \end{beamercolorbox}%
  }%
  \vskip0pt%
}

\begin{document}

\begin{frame}
\titlepage
\vspace{-0.45cm}
\begin{center}
\scriptsize
\renewcommand{\arraystretch}{1.08}
\begin{tabular}{@{}lll@{}}
\textbf{Name} & \textbf{Email} & \textbf{Matriculation} \\
\midrule
Love & love.-@stud.th-deg.de & 12306406 \\
Jamal Dassrath & jamal.dassrath@stud.th-deg.de & 22301035 \\
Erbakan Ahmad & erbakan.ahmad@stud.th-deg.de & 12306435 \\
Muhammad Ahtisham Bhatti & muhammad.bhatti2@stud.th-deg.de & 22301502 \\
\end{tabular}
\renewcommand{\arraystretch}{1.0}
\end{center}
\end{frame}

\begin{frame}{15 Minute Storyline}
\begin{columns}
\column{0.52\textwidth}
\begin{enumerate}
    \item Define the professor's project requirement.
    \item Introduce Oxford Pets and the two tasks.
    \item Explain preprocessing and reproducibility.
    \item Compare a scratch CNN with transfer learning.
    \item Read the generated results and limitations.
\end{enumerate}
\column{0.45\textwidth}
\IfFileExists{../outputs/figures/dataset_sample_grid.png}{
\includegraphics[width=\textwidth]{../outputs/figures/dataset_sample_grid.png}
}{}
\end{columns}
\end{frame}

\begin{frame}{Project Requirement Mapping}
\begin{center}
\begin{tabular}{p{0.38\textwidth}p{0.5\textwidth}}
\toprule
Requirement & Our implementation \\
\midrule
Two tasks on one dataset & 37-class breed classification and binary cat-vs-dog classification \\
Two deep learning approaches & Custom CNN from scratch and ResNet18 transfer learning \\
Course concepts & tensors, dataloaders, CNN layers, activations, cross-entropy, AdamW, validation, regularization \\
Professional deliverables & package code, notebook, report, slides, metrics, plots, checkpoints \\
\bottomrule
\end{tabular}
\end{center}
\end{frame}

\begin{frame}{Lecture Material We Used}
\begin{center}
\begin{tabular}{p{0.37\textwidth}p{0.52\textwidth}}
\toprule
Lecture topic & How it appears in our project \\
\midrule
Tensors and PyTorch & image batches shaped as \texttt{N x C x H x W} \\
Layers and activations & Conv2d, BatchNorm, ReLU, MaxPool, Linear \\
Loss and backpropagation & cross-entropy loss with autograd and optimizer steps \\
Optimizers & AdamW with learning rate and weight decay \\
Regularization & dropout, augmentation, validation split, early comparison \\
Convolution & local feature extraction for image classification \\
\bottomrule
\end{tabular}
\end{center}
\vspace{0.2cm}
\centering
\small We also checked the provided notebooks on tensors, training loops, datasets, optimizers, dropout, validation, and CNNs while building the project.
\end{frame}

\begin{frame}{Dataset: Oxford-IIIT Pets}
\IfFileExists{../outputs/figures/dataset_sample_grid.png}{
\centering
\includegraphics[height=0.78\textheight,width=0.95\linewidth,keepaspectratio]{../outputs/figures/dataset_sample_grid.png}
}{}
\end{frame}

\begin{frame}{Why This Dataset Is Useful}
\begin{columns}
\column{0.42\textwidth}
\begin{itemize}
    \item 37 pet categories.
    \item Roughly 200 images per category.
    \item Real-world variation in pose, background, lighting, and scale.
    \item Labels support both multi-class and binary classification.
\end{itemize}
\column{0.56\textwidth}
\IfFileExists{../outputs/figures/dataset_label_distribution.png}{
\includegraphics[width=\textwidth]{../outputs/figures/dataset_label_distribution.png}
}{}
\end{columns}
\end{frame}

\begin{frame}{Task Design}
\begin{columns}
\column{0.48\textwidth}
\textbf{Task 1: Breed classification}
\begin{itemize}
    \item 37 output classes.
    \item Harder task.
    \item Similar-looking breeds cause confusion.
\end{itemize}
\column{0.48\textwidth}
\textbf{Task 2: Species classification}
\begin{itemize}
    \item 2 output classes: cat and dog.
    \item Easier task.
    \item Useful sanity check for model learning.
\end{itemize}
\end{columns}
\end{frame}

\begin{frame}{Preprocessing and Augmentation}
\begin{columns}
\column{0.38\textwidth}
\begin{itemize}
    \item Resize and crop images.
    \item Random resized crop for training.
    \item Horizontal flip for augmentation.
    \item ImageNet normalization for transfer learning.
    \item Fixed seed for reproducibility.
\end{itemize}
\column{0.6\textwidth}
\IfFileExists{../outputs/figures/augmentation_preview.png}{
\includegraphics[width=\textwidth]{../outputs/figures/augmentation_preview.png}
}{}
\end{columns}
\end{frame}

\begin{frame}{Data Pipeline}
\begin{center}
\begin{tabular}{ll}
\toprule
Step & Implementation \\
\midrule
Download & \texttt{torchvision.datasets.OxfordIIITPet} \\
Splits & official trainval/test, plus validation split \\
Sampling & stratified smaller subsets for the laptop run \\
Batching & PyTorch \texttt{DataLoader} \\
Metrics & accuracy, macro F1, confusion matrix \\
Artifacts & CSV, JSON, PNG figures, model checkpoints \\
\bottomrule
\end{tabular}
\end{center}
\end{frame}

\begin{frame}{Model 1: CNN from Scratch}
\begin{columns}
\column{0.48\textwidth}
\begin{itemize}
    \item Learns all image features from Oxford Pets.
    \item Four convolutional blocks.
    \item Conv2d + BatchNorm + ReLU + MaxPool.
    \item Adaptive average pooling before classifier.
    \item Dropout and weight decay for regularization.
\end{itemize}
\column{0.48\textwidth}
\begin{center}
\begin{tabular}{ll}
\toprule
Layer idea & Course topic \\
\midrule
Convolution & local feature extraction \\
ReLU & non-linearity \\
BatchNorm & stable training \\
Dropout & regularization \\
Linear head & classification logits \\
\bottomrule
\end{tabular}
\end{center}
\end{columns}
\end{frame}

\begin{frame}{Model 2: Transfer Learning}
\begin{columns}
\column{0.48\textwidth}
\begin{itemize}
    \item ResNet18 pretrained on ImageNet.
    \item Replace final classifier for 37 or 2 classes.
    \item Train classifier head with frozen backbone.
    \item Fine-tune final ResNet block.
\end{itemize}
\column{0.48\textwidth}
\textbf{Why it should help}
\begin{itemize}
    \item Early CNN layers already detect edges and textures.
    \item Pet images are natural images, like ImageNet.
    \item Needs fewer examples than a scratch model.
\end{itemize}
\end{columns}
\end{frame}

\begin{frame}{Training Setup}
\begin{center}
\begin{tabular}{ll}
\toprule
Setting & Laptop run \\
\midrule
Image size & $96 \times 96$ \\
Optimizer & AdamW \\
Loss & Cross-entropy \\
Scratch training & 3 epochs \\
Transfer training & 3 frozen-head epochs + 1 fine-tuning epoch \\
Breed test subset & 370 stratified images \\
Species test subset & 400 stratified images \\
\bottomrule
\end{tabular}
\end{center}
\vspace{0.3cm}
\centering
\texttt{make train-laptop-final \quad make evaluate \quad make visuals}
\end{frame}

\begin{frame}{Overall Result}
\IfFileExists{../outputs/figures/model_comparison_summary.png}{
\centering
\includegraphics[width=0.9\linewidth]{../outputs/figures/model_comparison_summary.png}
}{}
\end{frame}

\begin{frame}{Final Metrics}
\begin{center}
\begin{tabular}{llrr}
\toprule
Task & Model & Accuracy & Macro F1 \\
\midrule
Breed & Scratch CNN & 0.049 & 0.014 \\
Breed & ResNet18 transfer & 0.308 & 0.268 \\
Species & Scratch CNN & 0.685 & 0.529 \\
Species & ResNet18 transfer & 0.898 & 0.877 \\
\bottomrule
\end{tabular}
\end{center}
\vspace{0.2cm}
\begin{itemize}
    \item Transfer learning improves both tasks.
    \item Species classification is much easier than breed classification.
    \item Macro F1 is useful because it reflects class-wise behavior, not only total accuracy.
\end{itemize}
\end{frame}

\begin{frame}{Breed Training Curves}
\begin{columns}
\column{0.5\textwidth}
\IfFileExists{../outputs/figures/breed_scratch_history.png}{
\includegraphics[width=\textwidth]{../outputs/figures/breed_scratch_history.png}
}{}
\column{0.5\textwidth}
\IfFileExists{../outputs/figures/breed_transfer_history.png}{
\includegraphics[width=\textwidth]{../outputs/figures/breed_transfer_history.png}
}{}
\end{columns}
\end{frame}

\begin{frame}{Species Training Curves}
\begin{columns}
\column{0.5\textwidth}
\IfFileExists{../outputs/figures/species_scratch_history.png}{
\includegraphics[width=\textwidth]{../outputs/figures/species_scratch_history.png}
}{}
\column{0.5\textwidth}
\IfFileExists{../outputs/figures/species_transfer_history.png}{
\includegraphics[width=\textwidth]{../outputs/figures/species_transfer_history.png}
}{}
\end{columns}
\end{frame}

\begin{frame}{Species Confusion Matrices}
\begin{columns}
\column{0.5\textwidth}
\IfFileExists{../outputs/figures/species_scratch_confusion_matrix.png}{
\includegraphics[width=\textwidth]{../outputs/figures/species_scratch_confusion_matrix.png}
}{}
\column{0.5\textwidth}
\IfFileExists{../outputs/figures/species_transfer_confusion_matrix.png}{
\includegraphics[width=\textwidth]{../outputs/figures/species_transfer_confusion_matrix.png}
}{}
\end{columns}
\end{frame}

\begin{frame}{Breed Confusion Matrix}
\begin{columns}
\column{0.54\textwidth}
\IfFileExists{../outputs/figures/breed_transfer_confusion_matrix.png}{
\includegraphics[width=\textwidth]{../outputs/figures/breed_transfer_confusion_matrix.png}
}{}
\column{0.42\textwidth}
\begin{itemize}
    \item Breed classification remains difficult.
    \item Some breeds share color, shape, fur, and facial structure.
    \item More epochs and larger images would help.
    \item Transfer learning still gives a clear improvement over scratch.
\end{itemize}
\end{columns}
\end{frame}

\begin{frame}{Prediction Examples}
\IfFileExists{../outputs/figures/species_transfer_prediction_slide.png}{
\centering
\includegraphics[height=0.82\textheight,width=0.96\linewidth,keepaspectratio]{../outputs/figures/species_transfer_prediction_slide.png}
}{}
\end{frame}

\begin{frame}{Interpretation}
\begin{itemize}
    \item The scratch CNN is valuable for learning the mechanics of deep learning.
    \item The transfer model is stronger because it starts from general visual features.
    \item The biggest performance jump is visible in breed classification, where learning features from scratch is hardest.
    \item Species classification validates that both pipelines can learn meaningful visual patterns.
\end{itemize}
\end{frame}

\begin{frame}{Limitations and Next Steps}
\begin{columns}
\column{0.48\textwidth}
\textbf{Limitations}
\begin{itemize}
    \item Our laptop run uses stratified subsets.
    \item Image size is reduced to $96 \times 96$.
    \item Short training favors transfer learning.
\end{itemize}
\column{0.48\textwidth}
\textbf{Next steps}
\begin{itemize}
    \item Run full dataset at $160 \times 160$.
    \item Train longer and tune learning rates.
    \item Add Grad-CAM interpretation.
    \item Try segmentation using trimaps.
\end{itemize}
\end{columns}
\end{frame}

\begin{frame}{Team Workflow}
\begin{center}
\begin{tabular}{ll}
\toprule
Student & Responsibility \\
\midrule
Love & dataset, EDA, preprocessing figures \\
Jamal Dassrath & scratch CNN implementation and experiments \\
Erbakan Ahmad & transfer learning implementation and experiments \\
Muhammad Ahtisham Bhatti & report, slides, final comparison, citations \\
\bottomrule
\end{tabular}
\end{center}
\vspace{0.3cm}
\begin{itemize}
    \item Shared commands: \texttt{make setup}, \texttt{make data}, \texttt{make train-laptop-final}.
    \item Result files can be regenerated and are ignored by Git until needed for submission.
\end{itemize}
\end{frame}

\begin{frame}{Conclusion}
\begin{itemize}
    \item The project satisfies the two-task and two-model assignment requirement.
    \item Transfer learning performs best on both breed and species classification.
    \item The scratch CNN remains important because it shows the course fundamentals.
    \item The project is packaged so every teammate can rerun, inspect, and extend it.
\end{itemize}
\end{frame}

\begin{frame}{Sources}
\begin{itemize}
    \item Oxford-IIIT Pet Dataset: \url{https://www.robots.ox.ac.uk/~vgg/data/pets/}
    \item Torchvision OxfordIIITPet documentation: \url{https://docs.pytorch.org/vision/main/generated/torchvision.datasets.OxfordIIITPet.html}
    \item PyTorch documentation: \url{https://pytorch.org/get-started/locally/}
    \item He et al., Deep Residual Learning for Image Recognition, CVPR 2016.
\end{itemize}
\end{frame}

\end{document}
